\section*{Bab \ref{ch:g12:metamorphic-rocks} --- Batuan Metamorf}

\begin{solution}{exo:g12:metamorphic-rocks:1}
(a) \omterm{def:g12:metamorphic-rocks:polymorph}{polimorf}; (b) \omterm{def:g12:metamorphic-rocks:index-mineral}{isograd}; (c) \omterm{def:g12:metamorphic-rocks:contact-regional}{aureol}; (d) \omterm{def:g12:metamorphic-rocks:foliation}{foliasi}; (e) \omterm{def:g12:metamorphic-rocks:lithostatic}{tekanan litostatik}.
\end{solution}

\begin{solution}{exo:g12:metamorphic-rocks:2}
(1) marmer; (2) kuarsit; (3) batu sabak; (4) gneis; (5) hornfels.
\end{solution}

\begin{solution}{exo:g12:metamorphic-rocks:3}
$P = \rho g z$: pada \qty{10}{km}, $2830 \times 9.81 \times \num{10000} \approx
\qty{2.78e8}{Pa} = \qty{0.278}{GPa}$; pada \qty{20}{km}, dua kali lipatnya,
\qty{0.555}{GPa}.
\end{solution}

\begin{solution}{exo:g12:metamorphic-rocks:4}
(1) Kianit (garis kianit--andalusit berada di dekat \qty{0.25}{GPa} pada
\qty{400}{\celsius}). (2) Andalusit. (3) Silimanit.
\end{solution}

\begin{solution}{exo:g12:metamorphic-rocks:5}
(1) Batu sabak, filit, sekis, gneis. (2) Pada sekis. (3) Gneis.
\end{solution}

\begin{solution}{exo:g12:metamorphic-rocks:6}
(1) $2830 \times 9.81 \times \num{35000} \approx \qty{9.72e8}{Pa}$, yaitu
$\num{9.72e8}/101325 \approx \num{9600}$ atmosfer. (2) $1025 \times 9.81
\times 4000 \approx \qty{4.02e7}{Pa} = \qty{40.2}{MPa}$. (3) $972/40.2
\approx 24$ kali.
\end{solution}

\begin{solution}{exo:g12:metamorphic-rocks:7}
(1) $z = \num{0.375e9}/(2830 \times 9.81) \approx \qty{13.5}{km}$.
(2) $(504 - 10)/13.5 \approx \qty{36.6}{\celsius}$ per kilometer. (3) Tidak:
batuan itu harus berada tepat pada satu titik diagram; batuan seperti itu
jarang.
\end{solution}

\begin{solution}{exo:g12:metamorphic-rocks:8}
(1) $23/3 \approx \qty{8}{mm}$. (2) Sekitar 8 juta tahun. (3) Batuan itu
tinggal di kedalaman, pada suhu tinggi, selama jutaan tahun.
\end{solution}

\begin{solution}{exo:g12:metamorphic-rocks:9}
(1) Empat: \omterm{def:g12:metamorphic-rocks:index-mineral}{isograd} garnet, staurolit, kianit, dan silimanit, dalam urutan itu.
(2) Ke arah barat. (3) Di ujung barat, di zona silimanit. Batuan-batuannya
telah melintasi garis kianit--silimanit: terkubur sekitar \qty{20}{km}
(\qty{0.56}{GPa}), batuan itu lebih panas daripada sekitar \qty{590}{\celsius}.
\end{solution}

\begin{solution}{exo:g12:metamorphic-rocks:10}
(1) Lumpur menjadi hornfels, batu kapur menjadi marmer, batu pasir menjadi
kuarsit. (2) Batuan-batuan itu terpanaskan tanpa terjepit: tidak ada tekanan
berarah yang menjajarkan mineral-mineralnya. (3) Batuan di sekitarnya sudah
panas, sehingga magma itu hanya sedikit lagi memanaskannya.
\end{solution}

\begin{solution}{exo:g12:metamorphic-rocks:11}
(1) Melintang terhadap dorongan: bidang pipihnya menghadap ke timur dan barat,
sehingga serpihan-serpihan itu berdiri pada bidang-bidang utara--selatan.
(2) Tegak, membentang utara--selatan. (3) Butir kalsit dan kuarsa tidak pipih
dan tidak panjang, sehingga walaupun mengkristal ulang di bawah jepitan, tidak
ada yang dapat berjajar.
\end{solution}

\begin{solution}{exo:g12:metamorphic-rocks:12}
(1) $z = (T - 10)/30$, sehingga $P = 0.0278\,(T - 10)/30 \approx
0.000927\,(T - 10)$. (2) $0.000927\,(T - 10) = 0.577 + 0.00209\,(T - 600)$
memberikan $0.001163\,T \approx 0.668$, sehingga $T \approx \qty{574}{\celsius}$, $P
\approx 0.000927 \times 564 \approx \qty{0.523}{GPa}$ dan $z = (574 - 10)/30
\approx \qty{18.8}{km}$. (3) Kianit di atas kedalaman itu, silimanit di
bawahnya.
\end{solution}

\begin{solution}{exo:g12:metamorphic-rocks:13}
(1) $10 + 10 \times 35 = \qty{360}{\celsius}$. (2) $10 + 30 \times 35 =
\qty{1060}{\celsius}$. (3) Tekanan pada \qty{35}{km} sekitar
\qty{0.97}{GPa}; pada \qty{360}{\celsius} titik itu terletak jauh di dalam
medan kianit. Titik kedua terletak jauh di sebelah kanan diagram, di medan
silimanit, pada suhu ketika batuan kerak mulai meleleh. (4) Keduanya memerlukan
tekanan tinggi pada suhu rendah, gabungan yang hanya disediakan oleh kerak
dingin yang terbawa turun di \omterm{def:g11:subduction-zones:subduction}{zona subduksi}.
\end{solution}

\begin{solution}{exo:g12:metamorphic-rocks:14}
(1) Kianit $1000/3.61 \approx \qty{277}{cm^3}$; silimanit $1000/3.24
\approx \qty{309}{cm^3}$; andalusit $1000/3.15 \approx \qty{317}{cm^3}$.
(2) $(317 - 277)/317 \approx 0.127$, sekitar \qty{13}{\percent}. (3) Pada
tekanan tinggi, susunan atom yang paling rapat diuntungkan; kianit memuat
atom-atom yang sama dalam ruang yang paling kecil.
\end{solution}

\begin{solution}{exo:g12:metamorphic-rocks:15}
(1) \qty{0.555}{GPa} pada \qty{20}{km}, \qty{0.694}{GPa} pada \qty{25}{km}.
(2) $T = 600 + (P - 0.577)/0.00209$: sekitar \qty{590}{\celsius} pada
\qty{0.555}{GPa} dan \qty{656}{\celsius} pada \qty{0.694}{GPa}. Pada
\qty{700}{\celsius} batuan itu lebih panas daripada keduanya, di medan
silimanit: sesuai. (3) $(700 - 10)/20 \approx \qty{34.5}{\celsius}$ dan
$(700 - 10)/25 \approx \qty{27.6}{\celsius}$ per kilometer, dekat dengan
gradien yang umum. (4) $\qty{20}{km}$ sampai $\qty{25}{km}$ dalam 40 juta tahun:
0.50 sampai \qty{0.63}{km} per sejuta tahun, yaitu 0.50 sampai
\qty{0.63}{mm} setahun.
\end{solution}

\begin{solution}{pb:g12:metamorphic-rocks:1}
\textbf{1.} Batuan itu berfoliasi, dan mikanya cukup besar untuk dilihat.
Mika-mikanya, muskovit dan biotit, membentuk \omterm{def:g12:metamorphic-rocks:foliation}{skistositasnya}.

\textbf{2.} \ce{Al2SiO5}; mineral-mineral itu adalah \omterm{def:g12:metamorphic-rocks:polymorph}{polimorf}.

\textbf{3.} Volume satu kilogram keduanya berbanding $3.15/3.61 \approx
0.873$: kianit sekitar \qty{13}{\percent} lebih kecil. Tekanan menguntungkan
kianit.

\textbf{4.} Zona kianit. Dalam perjalanan turun, batuan itu melintasi \omterm{def:g12:metamorphic-rocks:index-mineral}{isograd}
biotit, garnet, staurolit, dan kianit.

\textbf{5.} Sekitar 8 juta tahun.

\textbf{6.} $P = 0.577 + 0.00209 \times 0 = \qty{0.577}{GPa}$.

\textbf{7.} \qty{0.70}{GPa} berada di atas \qty{0.577}{GPa}, di sisi kianit
dari garis kianit--silimanit; titik itu juga jauh di atas medan andalusit, yang
pada \qty{600}{\celsius} terletak di bawah sekitar \qty{0.23}{GPa}.

\textbf{8.} \qty{0.577}{GPa}.

\textbf{9.} $0.70 = 0.577 + 0.00209\,(T - 600)$, sehingga $T = 600 +
0.123/0.00209 \approx \qty{659}{\celsius}$.

\textbf{10.} \qty{7.0}{kbar}; $\num{7.0e8}/101325 \approx \num{6900}$
atmosfer.

\textbf{11.} $\rho g = 2830 \times 9.81 \approx \qty{27800}{Pa}$ per meter,
\qty{27.8}{MPa} per kilometer.

\textbf{12.} $z = \num{0.70e9}/\num{27762} \approx \qty{25200}{m}$, sekitar
\qty{25.2}{km}.

\textbf{13.} $\num{0.577e9}/\num{27762} \approx \qty{20.8}{km}$.

\textbf{14.} $\num{0.70e9}/(2700 \times 9.81) \approx \qty{26.4}{km}$, sekitar
\qty{5}{\percent} lebih dalam ($26.4/25.2 \approx 1.048$).

\textbf{15.} Di dalam kerak: \qty{25}{km} bahkan kurang dari tebal kerak yang
normal.

\textbf{16.} $(600 - 10)/25.2 \approx \qty{23.4}{\celsius}$ per kilometer.

\textbf{17.} Gradien itu jauh lebih hangat daripada \omterm{def:g11:subduction-zones:subduction}{zona subduksi} (di bawah
\qty{10}{\celsius} per kilometer) dan sedikit lebih dingin daripada kerak yang
umum (\qty{30}{\celsius}), jauh dari daerah gunung api (40 sampai
\qty{50}{\celsius}): batuan yang terkubur dalam di kerak sabuk pegunungan yang
menebal, oleh \omterm{def:g12:metamorphic-rocks:contact-regional}{metamorfisme regional}.

\textbf{18.} $10 + 30 \times 25.2 \approx \qty{766}{\celsius}$. Garis
kianit--silimanit, bila diteruskan, terletak pada $0.577 + 0.00209 \times 166
\approx \qty{0.92}{GPa}$ pada suhu itu; \qty{0.70}{GPa} berada di bawahnya,
di medan silimanit. \omterm{def:g11:lithosphere-asthenosphere:geotherm}{Geoterm} yang umum akan menghasilkan silimanit.

\textbf{19.} $\qty{25.2}{km} = \qty{2.52e7}{mm}$; dengan laju \qty{0.5}{mm}
setahun, sekitar 50 juta tahun.

\textbf{20.} Ketika batuan itu naik, batuan itu mendingin, sehingga reaksinya
melambat, dan air yang diperlukan reaksi kebalikannya telah terusir dalam
perjalanan turun.

\textbf{21.} Sekis kianit itu mengkristal sekitar \textbf{\qty{25}{km} di
bawah permukaan}, pada sekitar \qty{600}{\celsius}.
\end{solution}
